% !TeX root = ../main.tex
\exercise{۲.۲۰}\label{sec:ex20}

\subsection*{صورت مسئله و محدوده مدل}

هدف، تعیین ولتاژ خروجی چهار مدار شکل \exnum{۲.۳۵} کتاب با صرف‌نظر از \lr{Body Effect} است \cite[صفحه چاپی ۹۷]{weste}. اتصال‌ها مستقیماً با شکل کتاب تطبیق داده شده‌اند. خط افقی تغذیه در شکل نشان‌دهنده \(V_{\mathrm{DD}}\)، نماد زمین نشان‌دهنده صفر ولت و حباب \lr{Gate} مشخص‌کننده \lr{pMOS} است.

مطابق تحلیل \lr{Threshold Drop} کتاب، از مدل \lr{Cutoff} با آستانه ثابت، بدون جریان بار \lr{DC} و بدون \lr{Leakage} استفاده می‌کنیم \cite[بخش \exnum{۲.۵.۴}]{weste} و \cite[اسلاید ۱۲]{dcslides}. منظور از خروجی‌های دارای افت، حد شارژ یا دشارژ در این مدل ساده است. این پاسخ بر رسیدن گره‌ها از سمت قابل‌هدایت به حد موردنظر تکیه دارد؛ برای گره شناور با شرایط اولیه دلخواه، الزاماً یک نقطه \lr{DC} یکتا وجود ندارد. کتاب نیز در پانوشت صفحه چاپی ۹۲، تغییر بسیار آهسته خروجی بر اثر \lr{Subthreshold Leakage} را یادآوری می‌کند.

برای نتایج متعارف زیر، تغذیه باید امکان روشن‌شدن ترانزیستورهای مربوط را بدهد: در بخش \lr{(b)}، \(V_{\mathrm{DD}}>2|V_{tp}|\) و در بخش \lr{(c)}، \(V_{\mathrm{DD}}>|V_{tp}|+V_{tn}\) فرض می‌شود. این‌ها شرط اعتبار تحلیل‌اند و داده عددی تازه‌ای برای مسئله نیستند.

\subsection*{قواعد انتقال ولتاژ}

در \lr{nMOS}، پایانه دارای ولتاژ کمتر نقش \lr{Source} را دارد. ترانزیستور با \(V_{gs}>V_{tn}\) هدایت می‌کند. بنابراین با \lr{Gate} در \(V_{\mathrm{DD}}\)، صفر را بدون افت آستانه عبور می‌دهد و هنگام شارژ یک گره از ولتاژ پایین، حد بالای آن برابر است با:

\nopagebreak[3]
\begin{equation}\label{eq:ex20-1}
V_{\mathrm{out,lim}}=V_g-V_{tn}=V_{\mathrm{DD}}-V_{tn}.
\end{equation}

در \lr{pMOS}، پایانه دارای ولتاژ بیشتر نقش \lr{Source} را دارد. شرط هدایت \(V_{sg}>|V_{tp}|\) است. این ترانزیستور با \lr{Gate} صفر، سطح \lr{HIGH} را بدون افت آستانه عبور می‌دهد؛ هنگام دشارژ یک گره به سمت زمین، حد پایین هدایت برابر است با:

\nopagebreak[3]
\begin{equation}\label{eq:ex20-2}
V_{\mathrm{out,lim}}=V_g+|V_{tp}|.
\end{equation}

در مرز آستانه، جریان مدل ایدئال به صفر می‌رسد. ازاین‌رو، برابری در دو رابطه بالا حد هدایت را مشخص می‌کند \cite[صفحات چاپی ۹۲ و ۹۳]{weste}.

\subsection*{تحلیل چهار مدار}

\subsubsection*{مدار \lr{(a)}: دو \lr{nMOS} در مسیر زمین}

ورودی ترانزیستور نخست زمین است و \lr{Gate} هر دو \lr{nMOS} به \(V_{\mathrm{DD}}\) متصل است. ترانزیستور نخست گره میانی را به صفر می‌رساند و ترانزیستور دوم نیز همان صفر را به خروجی منتقل می‌کند. با \(V_{\mathrm{DD}}>V_{tn}\)، مسیر دشارژ تا صفر هدایت دارد؛ در نتیجه:

\nopagebreak[3]
\begin{equation}\label{eq:ex20-3}
\boxed{V_{\mathrm{out}}=\qty{0}{\volt}}.
\end{equation}

سری‌شدن این دو ترانزیستور افت آستانه‌ای به سطح \lr{LOW} اضافه نمی‌کند؛ البته بر جریان و زمان دشارژ اثر می‌گذارد.

\subsubsection*{مدار \lr{(b)}: خروجی \lr{pMOS} نخست متصل به \lr{Gate} ترانزیستور دوم}

دو ترانزیستور در این مدار در یک مسیر \lr{Source–Drain} سری قرار ندارند. پایانه چپ و \lr{Gate} ترانزیستور نخست زمین‌اند؛ پایانه راست آن، گره \(X\)، به \lr{Gate} ترانزیستور دوم متصل است. پایانه چپ ترانزیستور دوم نیز زمین است و پایانه راست آن خروجی مدار را تشکیل می‌دهد.

هنگام دشارژ گره \(X\) از ولتاژ بالاتر، شرط هدایت ترانزیستور نخست \(V_X>|V_{tp}|\) است؛ بنابراین حد دشارژ برابر است با:

\nopagebreak[3]
\begin{equation}\label{eq:ex20-4}
V_{X,\mathrm{lim}}=|V_{tp}|.
\end{equation}

اگر خروجی ترانزیستور دوم از ولتاژ بالاتر دشارژ شود، شرط هدایت آن \(V_{\mathrm{out}}-V_X>|V_{tp}|\) خواهد بود. با رسیدن گره \(X\) به حد یادشده، حد دشارژ خروجی چنین است:

\nopagebreak[3]
\begin{equation}\label{eq:ex20-5}
\boxed{V_{\mathrm{out,lim}}=V_X+|V_{tp}|=2|V_{tp}|}.
\end{equation}

افت اضافی از کنترل \lr{Gate} دوم با یک سطح \lr{LOW} تضعیف‌شده ناشی می‌شود. اگر گره \(X\) از ابتدا زیر حد آستانه باشد، مدل بدون \lr{Leakage} آن را به \(|V_{tp}|\) شارژ نمی‌کند؛ بنابراین مقدار فوق را نباید مستقل از شرایط اولیه دلخواه، یک خروجی \lr{DC} تضمین‌شده دانست.

\subsubsection*{مدار \lr{(c)}: \lr{pMOS} و \lr{nMOS} در مسیر زمین}

\lr{Gate} ترانزیستور \lr{pMOS} نخست صفر است و ورودی آن نیز به زمین متصل است. گره میانی \(X\) در دشارژ از ولتاژ بالاتر تا حد زیر پایین می‌آید:

\nopagebreak[3]
\begin{equation}\label{eq:ex20-6}
V_{X,\mathrm{lim}}=|V_{tp}|.
\end{equation}

\lr{Gate} ترانزیستور \lr{nMOS} دوم به \(V_{\mathrm{DD}}\) متصل است. با شرط \(V_{\mathrm{DD}}-|V_{tp}|>V_{tn}\)، این ترانزیستور در حوالی حد نهایی همچنان هدایت دارد و دو گره میانی و خروجی را هم‌ولتاژ می‌کند. پس:

\nopagebreak[3]
\begin{equation}\label{eq:ex20-7}
\boxed{V_{\mathrm{out,lim}}=|V_{tp}|}.
\end{equation}

در این مدار، \lr{nMOS} افت آستانه تازه‌ای به مقدار \lr{LOW} تضعیف‌شده اضافه نمی‌کند. در مدل بدون \lr{Leakage}، شرط اولیه گره‌ها همچنان در تعبیر حد دشارژ اهمیت دارد.

\subsubsection*{مدار \lr{(d)}: \lr{pMOS} و \lr{nMOS} در مسیر تغذیه}

ورودی \lr{pMOS} برابر \(V_{\mathrm{DD}}\) و \lr{Gate} آن صفر است. با \(V_{\mathrm{DD}}>|V_{tp}|\)، گره میانی به \(V_{\mathrm{DD}}\) شارژ می‌شود. \lr{nMOS} دوم نیز \lr{Gate} با ولتاژ \(V_{\mathrm{DD}}\) دارد؛ هنگام شارژ خروجی از سطح پایین، با افزایش خروجی، \(V_{gs}\) کاهش می‌یابد. در مرز \(V_{gs}=V_{tn}\)، جریان مدل به صفر می‌رسد و داریم:

\nopagebreak[3]
\begin{equation}\label{eq:ex20-8}
\boxed{V_{\mathrm{out,lim}}=V_{\mathrm{DD}}-V_{tn}}.
\end{equation}

در نتیجه، انتقال \lr{HIGH} توسط \lr{pMOS} افتی ایجاد نمی‌کند و افت نهایی از \lr{nMOS} ناشی می‌شود.

\subsection*{جمع‌بندی و تفسیر}

جدول~\ref{tab:pass} پاسخ متعارف مسئله در مدل \lr{Threshold Drop} و با شرایط شارژ یا دشارژ بیان‌شده را خلاصه می‌کند.

\begin{table}[htbp]\centering
\caption{خروجی شبکه‌های ترانزیستوری}
\label{tab:pass}
\begin{tabularx}{\linewidth}{ccX}\toprule
مدار & خروجی یا حد انتقال & علت\\
\midrule
\lr{(a)} & \(\qty{0}{\volt}\) & انتقال \lr{LOW} توسط دو \lr{nMOS}\\
\lr{(b)} & \(2|V_{tp}|\) & \lr{LOW} تضعیف‌شده در \lr{Gate} ترانزیستور \lr{pMOS} دوم\\
\lr{(c)} & \(|V_{tp}|\) & افت ناشی از \lr{pMOS} نخست و هدایت \lr{nMOS} دوم\\
\lr{(d)} & \(V_{\mathrm{DD}}-V_{tn}\) & محدودیت شارژ خروجی توسط \lr{nMOS}\\
\bottomrule\end{tabularx}\end{table}

نتیجه بخش \lr{(b)} با سری‌کردن ساده دو \lr{pMOS} در مسیر عبور سیگنال یکسان نیست. تفاوت اتصال \lr{Gate} و مسیر \lr{Source–Drain}، علت انباشته‌شدن افت در آن مدار است؛ این تمایز در بخش \exnum{۲.۵.۴} کتاب نیز توضیح داده شده است. این نتایج، پیش‌بینی مدل ساده تمرین‌اند و شامل محاسبه \lr{Leakage}، زمان نشست یا شبیه‌سازی یک فناوری واقعی نمی‌شوند.

